% ============================================================================
% ICSE Paper — Approach Section (v2b: functional knowledge structure)
% LLM-Driven Trace Link Recovery between Architecture Documents and Models
% ============================================================================

\section{Approach}\label{sec:approach}

\subsection{Overview}\label{sec:overview}

We propose an LLM-driven pipeline for recovering trace links between
Software Architecture Documents (SAD) and Software Architecture Models (SAM).
Given a natural-language architecture document consisting of sentences
$\{s_1, \ldots, s_n\}$ and a component model containing components
$\{c_1, \ldots, c_m\}$, our approach produces a set of trace links
$L \subseteq S \times C$, where each link $(s_i, c_j)$ asserts that
sentence~$s_i$ discusses component~$c_j$.

\paragraph{Challenge.}
Architecture documents refer to components in diverse ways---by
canonical name, abbreviation, functional description, trailing
word of a compound name (e.g., ``Dispatcher'' for
\texttt{EventDispatcher}), or implicitly via pronouns---and many
component names are inherently ambiguous: common-English words
like \texttt{Core} or components named after technologies
appear frequently in non-referential contexts.
These two dimensions---\emph{reference diversity} and \emph{name
ambiguity}---interact: capturing diverse reference types demands
broad extraction, but ambiguous names make broad extraction
error-prone.

Existing approaches address only part of this problem.
Lexical NLP pipelines such as SWATTR~\cite{TODO:swattr} rely on
string similarity and cannot resolve abstraction mismatches
(``order processing subsystem'' $\to$ \texttt{OrderCore}),
discourse dependencies (pronouns referring to components named
sentences earlier), or generic-word collisions.
LLM-based classifiers such as LISSA~\cite{TODO:lissa} add
semantic understanding but lack structured knowledge of
architectural naming conventions---no alias tables, no
ambiguity-aware classification, no coreference resolution---treating
each sentence independently.

\paragraph{Key insight.}
Reliable recovery requires combining LLM reasoning with
structured architectural knowledge acquired before linking begins.
Different reference types need different resolution strategies,
but these strategies can operate independently on a shared
knowledge base; and a convention-aware filter grounded in
architectural naming patterns catches false positives that
neither lexical matching nor unguided LLM classification can
identify.
This decomposes into three sub-problems, each handled by one
pipeline layer:
resolving the \emph{knowledge gap}---which names are ambiguous,
what aliases exist, and which terms denote component
internals (\S\ref{sec:knowledge});
recovering links via \emph{per-type strategies} that consume
this knowledge (\S\ref{sec:recovery}); and
consolidating candidates and \emph{filtering residual noise}
using the acquired architectural knowledge
(\S\ref{sec:merge}).

\paragraph{Design principles.}
Three principles guide the pipeline design:
\begin{itemize}
    \item \emph{Multi-pass consensus:}
      Precision-critical decisions use two independent LLM
      passes---with task-specific agreement rules---to suppress
      hallucinations and LLM non-determinism (\S\ref{sec:seed},
      \S\ref{sec:entity}, \S\ref{sec:exclusion}).
    \item \emph{Citation-grounded verification:}
      Rather than trusting LLM judgments as opaque decisions, we
      require the LLM to produce verifiable evidence (antecedent
      citations, exact substrings) that is checked deterministically
      in code (\S\ref{sec:coref}, \S\ref{sec:entity}).
    \item \emph{Benchmark-clean prompts:}
      All few-shot examples use abstract placeholders or safe
      textbook domains; no benchmark names or gold-standard
      information appears in any prompt.
\end{itemize}

\paragraph{Architecture.}
The pipeline is organized as a staged DAG
(Figure~\ref{fig:pipeline}) with three layers:

\begin{enumerate}
    \item \textbf{Knowledge acquisition} (Layer~1): Analyses
      extract three categories of structured knowledge---component
      name properties, an alias table, and exclusion terms---from the
      raw inputs.
      Some analyses are mutually independent (sub-tier~1a); others
      depend on earlier results and run in a second sub-tier (1b).
    \item \textbf{Link recovery} (Layer~2): Four complementary
      strategies produce candidate links, each targeting a distinct
      reference type (\S\ref{sec:recovery}). Their outputs converge
      only at the merge point.
    \item \textbf{Link consolidation} (Layer~3): A single fan-in
      stage merges all candidates, applies convention-aware boundary
      filtering to \emph{all} links regardless of provenance, and
      removes residual false positives via a deterministic evidence
      filter.
\end{enumerate}

\noindent Explicit reference extraction is logically a Layer~2
strategy (it produces links, not knowledge), but because it requires
no Layer~1 knowledge, it runs concurrently with Layer~1 to amortize
its LLM cost against the knowledge phases.

% ---------------------------------------------------------------------------
% ============================================================================
% TikZ figure for the pipeline DAG
% Requires: \usepackage{tikz}, \usetikzlibrary{positioning,arrows.meta,fit,backgrounds}
% ============================================================================

\begin{figure*}[t]
\centering
\begin{tikzpicture}[
    node distance=0.8cm and 0.6cm,
    phase/.style={draw, rounded corners, minimum width=2.2cm, minimum height=0.7cm,
                  font=\footnotesize, align=center},
    layer1/.style={phase, fill=blue!8},
    layer2/.style={phase, fill=orange!10},
    layer3/.style={phase, fill=green!8},
    arr/.style={-{Stealth[length=2mm]}, thick, gray!70},
    layerbox/.style={draw=gray!40, dashed, rounded corners=3pt, inner sep=6pt},
    layerlabel/.style={font=\scriptsize\sffamily, gray!60},
]

% Layer 1a: Independent analyses (top row)
\node[layer1] (model)   {Model\\Analysis};
\node[layer1, right=1.2cm of model] (doc) {Document\\Analysis};

% Layer 1b: Enrichment (needs 1a results)
\node[layer1, below=0.7cm of model, xshift=1.7cm] (enrich) {Partial-Ref.\\Refinement};

% Layer 1 internal dependencies
\draw[arr] (model.south east) -- (enrich.north west);
\draw[arr] (doc.south west) -- (enrich.north east);

% Layer 2: Link Recovery
\node[layer2, below=1.4cm of enrich, xshift=-2.0cm] (explicit) {Explicit\\Ref.\ Extr.};
\node[layer2, right=0.6cm of explicit] (contextual) {Contextual\\Ref.\ Disc.};
\node[layer2, right=0.6cm of contextual] (anaphoric) {Anaphoric\\Ref.\ Res.};
\node[layer2, right=0.6cm of anaphoric] (abbreviated) {Abbreviated\\Ref.\ Match};

% Knowledge flow arrows (Layer 1 -> Layer 2)
\draw[arr] (enrich.south west) -- (explicit.north east);
\draw[arr] (model.south) -- (contextual.north west);
\draw[arr] (enrich.south) -- (contextual.north);
\draw[arr] (enrich.south east) -- (anaphoric.north);
\draw[arr] (enrich.south east) -- (abbreviated.north west);

% Layer 3: Link Consolidation
\node[layer3, below=1.4cm of contextual, xshift=0.8cm] (merge) {Merge};

% Dedup dependency: abbreviated matching needs prior strategies' results
\draw[arr, dashed] (explicit.east) -- (abbreviated.west);

% Layer 2 -> Layer 3
\draw[arr] (explicit.south) -- (merge.north west);
\draw[arr] (contextual.south) -- (merge.north);
\draw[arr] (anaphoric.south) -- (merge.north);
\draw[arr] (abbreviated.south) -- (merge.north east);

% Output
\node[below=0.4cm of merge, font=\footnotesize\itshape] (out) {Trace Links $L$};
\draw[arr] (merge) -- (out);

% Layer labels
\node[layerlabel, left=0.3cm of model] {Layer 1a};
\node[layerlabel, left=0.3cm of enrich] {Layer 1b};
\node[layerlabel, left=0.3cm of explicit] {Layer 2};
\node[layerlabel, left=0.8cm of merge] {Layer 3};

% Layer boxes
\begin{scope}[on background layer]
    \node[layerbox, fit=(model)(doc)(enrich), label={[layerlabel]above:Knowledge Acquisition}] {};
    \node[layerbox, fit=(explicit)(abbreviated), label={[layerlabel]above left:Link Recovery}] {};
    \node[layerbox, fit=(merge), label={[layerlabel]above left:Link Consolidation}] {};
\end{scope}

\end{tikzpicture}
\caption{Pipeline architecture as a staged DAG. Layer~1 acquires knowledge
in two sub-tiers: model analysis and document analysis run
concurrently~(1a); partial-reference refinement depends on their
results~(1b).
Explicit reference \emph{extraction} has no knowledge dependencies
and starts concurrently with Layer~1 (shown here under Link Recovery
for grouping); its \emph{validation} uses the alias table once
Layer~1 completes.
Layer~2 runs four recovery strategies; abbreviated matching
depends on prior strategies for deduplication (dashed arrow).
Layer~3 deduplicates their output.}
\label{fig:pipeline}
\end{figure*}

% ---------------------------------------------------------------------------

\subsection{Knowledge Acquisition}\label{sec:knowledge}

Layer~1 resolves the three knowledge gaps identified in
\S\ref{sec:overview}---name ambiguity (\S\ref{sec:names}),
alternative names (\S\ref{sec:alias}), and internal-part
identification (\S\ref{sec:exclusion})---through targeted analyses
organized in two sub-tiers.
Sub-tier~1a extracts name properties and the core alias table
concurrently (each reads only raw inputs); sub-tier~1b enriches
the alias table and identifies exclusion terms once 1a completes.

\subsubsection{Component Name Properties}\label{sec:names}

Component names vary widely in specificity.
Constructed identifiers like \texttt{TaskScheduler} unambiguously name
architectural roles, while single common-English words like
\texttt{Core} or \texttt{Adapter} appear frequently in
non-referential contexts.
To distinguish them, each name is classified as \emph{architectural}
or \emph{ambiguous} via a hybrid of LLM classification and
deterministic structural guards: names with unambiguous structure
(multi-word, CamelCase, all-uppercase) are always classified as
architectural, ensuring LLM judgment is used only where structural
cues are absent.

The ambiguous/architectural distinction drives several downstream
decisions: which candidates require generic-mention detection during
validation (\S\ref{sec:entity}), which trailing words need
capitalized matching during alias enrichment (\S\ref{sec:alias}),
and which names receive stricter scrutiny in the boundary filter
(\S\ref{sec:boundary}).

Name analysis also identifies \emph{implementation relationships}
between components
(e.g., \texttt{BasicScheduler}~$\to$~\texttt{Scheduler})
by shared-substring matching.
Implementation variants are excluded from all downstream LLM prompts
(which see only the abstract parent) and deduplicated at merge time
via a parent-overlap guard (\S\ref{sec:dedup}).

\emph{Runs in sub-tier~1a (concurrent, reads only the model).}

\subsubsection{Alias Table}\label{sec:alias}

Architecture documents rarely refer to components exclusively by
their canonical names.
Authors abbreviate (\emph{AST} for \texttt{AbstractSyntaxTree}),
use trailing words as shorthand (\emph{Dispatcher} for
\texttt{EventDispatcher}), or introduce synonyms.
Without an \emph{alias table}---a mapping from these alternative
names to canonical components---any name-matching strategy will
systematically miss such references, capping recall.
The alias table feeds contextual discovery (\S\ref{sec:entity}),
abbreviated matching (\S\ref{sec:partial}), and the boundary filter
(\S\ref{sec:boundary}).

\paragraph{Core discovery (sub-tier~1a).}
Because aliases arise from different linguistic
conventions, discovery must handle each category with
distinct validation rules.
An LLM extracts three categories: abbreviations
(parenthetical definitions), synonyms (unambiguous alternative names),
and partial references (trailing words of multi-word names used
alone).
A calibrated few-shot judge validates each proposed mapping,
rejecting generic terms (\emph{system}, \emph{process}) while
auto-approving structurally unambiguous identifiers.

\paragraph{Partial-reference refinement (sub-tier~1b).}
Among the three alias types, partial references require the most
careful handling because many trailing words are ordinary English
words whose referential status depends on context.
Refinement consumes both the initial alias table (to avoid
duplicates) and the ambiguous-name set from name analysis (to
distinguish generic trailing words from architectural ones),
placing it in sub-tier~1b.

A statistical enrichment step adds partial mappings when a trailing
word appears in $\ge3$~sentences where the full component name is
absent, and no other component shares that trailing word;
trailing words classified as ambiguous require capitalized occurrences
to avoid coincidental overlap.
A subsequent per-partial LLM classification determines whether each
generic partial is typically used as a component reference
(\textsc{name}) or as an ordinary English word (\textsc{ordinary}).
\textsc{Ordinary}-classified partials remain in the alias table
for deterministic matching (\S\ref{sec:partial}) but are
excluded from the confirmed-alias context given to the boundary
filter (\S\ref{sec:boundary}), preventing it from treating
generic uses as validated references.

\subsubsection{Exclusion Terms}\label{sec:exclusion}

Architecture documents describe not only components but also their
internal parts---subprocess terms like ``garbage collector'' inside a
\texttt{MemoryManager}.
If left unrecognized, these terms create false positives in
anaphoric resolution, where a sentence about an internal mechanism
is misattributed to the enclosing component.

Two independent LLM passes identify subprocess terms via a
cross-validation debate: the first pass proposes candidates, the
second validates them, rejecting terms that may be valid component
references.
Only terms confirmed by both passes are retained.
Subprocess terms serve as a hard filter in anaphoric resolution
(\S\ref{sec:coref}), where sentences classified as subprocess
descriptions are deterministically excluded from coreference linking,
and as contextual hints in validation prompts (\S\ref{sec:entity}).

\emph{Runs in sub-tier~1b (needs the filtered component list from
name analysis).}

% ---------------------------------------------------------------------------

\subsection{Link Recovery}\label{sec:recovery}

No single strategy can capture all reference types:
canonical mentions are lexically recognizable but aliases require
knowledge, pronouns need discourse context, and abbreviated names
demand project-specific awareness.
With shared knowledge in place, four complementary strategies
each address one of these types:
canonical mentions (\S\ref{sec:seed}),
alias-based and contextual references (\S\ref{sec:entity}),
pronominal coreferences (\S\ref{sec:coref}), and
abbreviated names (\S\ref{sec:partial}).
Their outputs converge only at the Layer~3 merge point.

\subsubsection{Explicit Reference Extraction}\label{sec:seed}

References using the canonical component name can be captured
without alias knowledge, but even these require care: ambiguous
names risk false positives, and names appearing only inside
hierarchical paths must not be treated as component references.
A self-contained LLM-based extractor (\textbf{ILinker2})
provides this high-precision baseline by reading only the raw
sentences and component list, enabling concurrent execution with
Layer~1.

ILinker2 uses two complementary LLM passes over batched document
segments, each with a different prompt framing:
Pass~A (mention-framed) asks which components are explicitly
mentioned; Pass~B (actor-framed) asks which component is the subject
or primary actor.
Both passes reject names appearing only inside hierarchical
namespace paths and generic English words used in their ordinary
sense.

The merge applies \emph{asymmetric voting}:
exact-name matches from \emph{either} pass are accepted (union),
while non-exact matches require agreement from \emph{both} passes
(intersection).
Exact matches are high-confidence and reliably identified by a single
pass; non-exact matches risk hallucination and benefit from
independent agreement.
ILinker2 achieves high precision at the cost of recall;
the recall gap---alias-based references, functional descriptions,
pronominal coreferences---is addressed by the three knowledge-enriched
strategies.

\subsubsection{Contextual Reference Discovery}\label{sec:entity}

Where explicit extraction finds only references recognizable from
the component name alone, contextual discovery leverages the full
alias table to recover references that \emph{require} contextual
understanding.
It differs from explicit extraction in three ways:
(1)~knowledge-enriched prompts receive all aliases, enabling
recognition of abbreviations and partial names that the
knowledge-free extractor cannot match;
(2)~broader inclusion scope covers space-separated compound forms,
functional descriptions, known aliases, interaction participants
(``X sends data to Y'' references both X and Y), and
passive/prepositional mentions (``handled by X'');
(3)~same-prompt intersection runs the \emph{same} prompt
\emph{twice} independently, keeping only candidates found in both
passes.
Because the broader scope increases the risk of over-extraction,
same-prompt intersection is necessary to suppress spurious candidates
that a single pass might hallucinate---a risk that grows with
inclusion scope.

For components that receive no links from either extraction or
discovery, a \emph{targeted recovery} step issues single-component
prompts with all known aliases, recovering references missed by
broad-scope batch processing.

\paragraph{Validation.}
The broader extraction scope increases the risk of false
positives, requiring multi-step validation before candidates
are accepted:
\begin{enumerate}
    \item \emph{Generic-mention detection:}
      For candidates involving an ambiguous component name in lowercase
      (e.g., ``core'' rather than ``Core''), an LLM determines whether
      the usage refers to the architectural component or is an ordinary
      English word.
      The LLM sees anchor sentences where the full name appears for
      calibration, and classifies each candidate as \textsc{component}
      or \textsc{generic}; generic-classified candidates are rejected
      before further validation.
    \item \emph{Code-first auto-approval:}
      Candidates where a component name or alias appears verbatim
      in the sentence text (verified by word-boundary matching in
      code) are approved without LLM cost.
    \item \emph{Two-pass intersection:}
      Remaining candidates are validated by two independent LLM passes
      with complementary foci---one on \emph{actor role}, one on
      \emph{direct reference}---and accepted only if both approve.
\end{enumerate}

\subsubsection{Anaphoric Reference Resolution}\label{sec:coref}

Sentences containing pronouns (\emph{it}, \emph{they}, \emph{this},
\emph{these}, \emph{that}, \emph{those}) and definite descriptions
(\emph{the component}, \emph{the service}) carry valid trace
information but contain no component name, making them invisible to
the preceding strategies.
However, the component a pronoun refers to was typically named
explicitly in a nearby sentence;
if the LLM can identify this antecedent and the identification
can be verified deterministically, the resolution becomes trustworthy
without additional review.

The pipeline uses a unified \emph{cases-in-context} strategy:
each pronoun-bearing sentence is presented as an individual case with
a $\pm5$-sentence bidirectional context window, processed in batches
of~10.
The LLM resolves any pronoun that refers to an architecture component
and must cite an antecedent sentence along with the exact quote
containing the component name.
This citation is verified deterministically---the name (or alias)
must appear verbatim in the cited sentence, the antecedent must be
within 3~sentences (tighter than the $\pm5$ context window
to suppress long-distance errors), and sentences classified as subprocess
descriptions (\S\ref{sec:exclusion}) are excluded---making
hallucinated resolutions where the LLM fabricates a nonexistent
antecedent detectable and rejectable in code.

Because the verification is strict (verbatim match within
3~sentences), false positives are effectively eliminated at the
code level, making anaphoric links safe to accept without further
review.

\subsubsection{Abbreviated Reference Matching}\label{sec:partial}

Contextual discovery may miss sentences where a partial name
appears without surrounding semantic cues.
Determining whether a partial name like ``Server'' refers to the
component or is used generically requires project-specific context
that a single LLM prompt cannot reliably encode.

Instead, the pipeline deterministically scans all sentences for
word-boundary matches of each partial reference from the alias table,
producing links for pairs not already covered by preceding strategies.
A clean-mention filter ensures the partial name appears as a
standalone word---not inside a dotted path, a hyphenated compound,
or a qualified name.
This strategy uses no LLM calls, providing stable recall at zero
additional inference cost;
its output is evaluated individually by the boundary filter in
Layer~3 (\S\ref{sec:boundary}).

% ---------------------------------------------------------------------------

\subsection{Link Consolidation}\label{sec:merge}

Independent strategies with different precision profiles produce
three categories of noise:
overlapping candidates from independent discovery of the same
link (\S\ref{sec:dedup}),
component names used in non-architectural senses that neither
lexical matchers nor generic LLM classifiers can detect
(\S\ref{sec:boundary}), and
hallucinated references lacking textual grounding
(\S\ref{sec:evidence-filter}).
Layer~3 addresses each with a dedicated stage, leveraging the
architectural knowledge from Layer~1.

\subsubsection{Priority-Based Merge}\label{sec:dedup}

Because recovery strategies operate independently, they may
discover the same $(s_i, c_j)$ pair via different evidence.
These overlaps are deduplicated, retaining the highest-priority
source:
$\text{explicit} > \text{contextual} > \text{anaphoric} >
\text{abbreviated}$,
reflecting decreasing levels of direct textual evidence.

A \emph{parent-overlap guard} removes links to implementation
components when the abstract parent is already linked to the same
sentence (e.g., if both \texttt{Scheduler} and
\texttt{BasicScheduler} are linked, the implementation variant is
dropped), using the implementation relationships from name
analysis (\S\ref{sec:names}).

\subsubsection{Convention-Aware Boundary Filter}\label{sec:boundary}

Even after merge, the name-ambiguity challenge persists: some
links survive where a component name appears in a sentence but
is used in a non-architectural sense.
A structured LLM filter applies a reasoning guide to \emph{all}
links, regardless of provenance, checking two rejection patterns:
\begin{enumerate}
    \item \emph{Hierarchical name reference:}
      The component name appears only as a prefix in a dotted or
      qualified path (e.g., ``\texttt{X.utils} provides helper
      functions'' does not reference component~X).
      An exception preserves links when the sentence also names the
      component explicitly.
    \item \emph{Name confusion:}
      The name is used as a technology, methodology, or generic
      English word rather than as the architectural component.
      A critical exception preserves links when a component is
      \emph{named after} a technology---any sentence describing that
      technology's capabilities is about the component.
\end{enumerate}

\noindent If neither rejection pattern applies, the link is kept.
Structural guardrails ensure multi-word names and CamelCase
identifiers are never treated as generic words, and that interaction
sentences (``X sends data to Y'') produce links for all
participating components.

\subsubsection{Evidence Filter}\label{sec:evidence-filter}

Even after boundary filtering, contextual discovery may have
produced links where the LLM hallucinated a reference not
present in the sentence text.
A deterministic filter catches these by requiring textual
grounding---a link is retained only if \emph{any} of the
following holds:
\begin{enumerate}
    \item The component name or a known alias appears as a
      word-boundary match in the sentence text.
    \item The link originates from explicit extraction or anaphoric
      resolution, both of which already verified the reference at
      their respective extraction stages.
\end{enumerate}

\noindent Links satisfying neither condition are dropped.
This filter requires no LLM calls and provides a final
deterministic safety net against ungrounded LLM outputs.
